\cotefile{111}
\scanpage{2}\mspage{\handwritten{\uline{1}}}\margin{\handwritten{notes du \uline{10.9.68}}}

Tapis de Quillen.

I \uline{Relations entre catégories et ensembles semi-simplic\unsure{aux|iaux}}.

A toute catégorie $C$, on associe un ensemble semi-simpli\deleted{\unsure{v}}cial $S(C)$, trouvant \unsure{ninsi|ainsi} un foncteur pleinement fidèle \deleted{de Cat}
\[
S;\ (\mathrm{Cat})\handwritten{\longrightarrow}(\mathrm{Ssimpl})\ .
\]

\deleted{Si $C$ est associé à} Les systèmes locaux d'ensemble sur $SC$ correspondent aux foncteurs sur $C$ qui transforment toute flêche en isomorphisme (i.e. qu\deleted{\unsure{e}}i se factorisent par le groupoïde associé à $C$). Les $H^i$ sur $SC$ d'un tel système local ($H^{\unsure{o|0}}$ pour ens, $H^1$ pour groupes, $H^i$ quelconques pour groupes abéliens) s'interprètent en termes des foncteurs $\underset{\handwritten{\leftarrow}}{\mathrm{lim}}^{(i)}$ déri\deleted{\unsure{v}}gés de $\underset{\handwritten{\leftarrow}}{\mathrm{lim}}$, ou si on préfère, des $H^i($ du \uline{topos} $C^{\handwritten{\wedge}}$. On voit ainsi à quelle condition un foncteur $C\handwritten{\to}C'$ induit un\deleted{\unsure{e}} homotopisme $SC\handwritten{\to}SC'$ : En vertu du critère cohomologique de Artin-Mazur, il f et s que pour tout système de coéfficients $F'$ sur $C'$, l'homomorphisme naturel $\underset{C'}{\underset{\handwritten{\leftarrow}}{\mathrm{lim}}}^{(i)}F'\handwritten{\longrightarrow}\underset{C}{\underset{\handwritten{\leftarrow}}{\mathrm{lim}}}^{(i)}F$ soit un isomorphisme (pour les $i$ pour lesquels cela a un sens). \deleted{En termes}

A $C$ on peut asso\deleted{\unsure{v}}c\unsure{er|ier|ér} le topos $\handwritten{\widetilde{C}}$, qui varie de façon \uline{covariante} avec $C$. (NB le foncteur $C\handwritten{\mapsto\widetilde{C}}$ n'a plus rien de pleinement fidèle, semble-til ??) Les systèmes de coéfficients \inserted{ensemblistes} sur $C$ ($\overset{\mathrm{déf}}{=}$ les foncteurs \inserted{$C^{o}\handwritten{\to}\mathrm{Ens}$} transformant isomorphismes en isomorphismes) correspondent aux faisceaux localement constants i.e. les objets localement constants de $\handwritten{\widetilde{C}}$, définis intrinsèquement en \unsure{termes} de $\handwritten{\widetilde{C}}$. Ainsi, le fait pour un foncteur $F:C\handwritten{\to}C'$ d'induire un\deleted{e équivalence} homotopisme $S(C)\handwritten{\to}S(C')$ ne dépend que du morphisme de topos $\handwritten{\widetilde{F}}:\handwritten{\widetilde{C}}\handwritten{\to}\handwritten{\widetilde{C}}'$ induit, et signifie que pour tout faisceau localement constant $F'$ sur $C'$ \unsure{i.}e. sur $\handwritten{\widetilde{C}}'$, les applications induites $H^i(\handwritten{\widetilde{C}}',F')\handwritten{\longrightarrow}H^i(\handwritten{\widetilde{C}},\handwritten{\widetilde{F}}^{\handwritten{*}}(F'))$ sont des isomorphismes (pour les $i$ pour lesquels cela a un\deleted{\unsure{es|e}} sens).

% archive sheet number: 1
